\section{El perceptrón: el bloque de construcción básico del aprendizaje profundo}

\subsection{Propagación hacia adelante del perceptrón}

El \textbf{perceptrón} (Perceptron) es la unidad estructural básica de todas las redes neuronales, también llamada «neurona» (neuron). Su proceso de propagación hacia adelante consta de tres pasos:

\begin{enumerate}
\item \textbf{Suma ponderada}: multiplicar cada entrada $x_i$ por su peso correspondiente $w_i$ y luego sumar
\item \textbf{Suma del sesgo}: sumar el término de sesgo $w_0$ (bias)
\item \textbf{Activación no lineal}: obtener la salida mediante una función de activación no lineal $g(\cdot)$
\end{enumerate}

\begin{figure}[H]
\centering
\includegraphics[width=0.85\textwidth]{slides-images/slide-022.jpg}
\caption{Propagación hacia adelante del perceptrón (con término de sesgo)\protect\footnotemark}
\end{figure}
\footnotetext{Diapositivas, página 22.}

Su expresión matemática es:

$$\hat{y} = g\left(w_0 + \sum_{i=1}^{m} x_i w_i\right) = g(w_0 + \mathbf{X}^T \mathbf{W})$$

donde:
\begin{itemize}
\item $\mathbf{X} = [x_1, x_2, \ldots, x_m]^T$: vector de entrada
\item $\mathbf{W} = [w_1, w_2, \ldots, w_m]^T$: vector de pesos
\item $w_0$: término de sesgo (bias), que permite desplazar la función de activación a lo largo del eje $z$
\item $g(\cdot)$: función de activación no lineal
\item $\hat{y}$: salida predicha
\end{itemize}

\begin{importantbox}{La función del término de sesgo}
El término de sesgo $w_0$ es un parámetro indispensable en el perceptrón. Permite desplazar la función de activación con independencia de los valores de entrada, de modo que la neurona pueda producir una salida distinta de cero incluso cuando todas las entradas son cero. Un perceptrón sin sesgo solo puede aprender fronteras de decisión que pasan por el origen.
\end{importantbox}

\subsection{Funciones de activación}

La función central de la función de activación $g(\cdot)$ es \textbf{introducir no linealidad en la red}. Entre las funciones de activación más comunes están:

\begin{figure}[H]
\centering
\includegraphics[width=0.95\textwidth]{slides-images/slide-025.jpg}
\caption{Tres funciones de activación comunes: Sigmoid, Tanh y ReLU\protect\footnotemark}
\end{figure}
\footnotetext{Diapositivas, página 25.}

\begin{table}[H]
\centering
\begin{tabular}{llll}
\toprule
\textbf{Función de activación} & \textbf{Fórmula} & \textbf{Rango de salida} & \textbf{Características} \\
\midrule
Sigmoid & $\sigma(z) = \frac{1}{1+e^{-z}}$ & $(0, 1)$ & Adecuada para salidas de probabilidad \\
Tanh & $g(z) = \frac{e^z - e^{-z}}{e^z + e^{-z}}$ & $(-1, 1)$ & Centrada en cero \\
ReLU & $g(z) = \max(0, z)$ & ${[}0, +\infty)$ & Computacionalmente eficiente, la más usada \\
\bottomrule
\end{tabular}
\caption{Comparación de funciones de activación comunes}
\end{table}

\begin{warningbox}{¿Por qué es indispensable usar una función de activación no lineal?}
Si no se usa una función de activación no lineal (es decir, si se usa una activación lineal $g(z)=z$), entonces, sin importar cuántas capas tenga la red, la salida de toda la red seguirá siendo solo una combinación lineal de la entrada. Esto significa que una red de múltiples capas se degrada a un modelo lineal de una sola capa, incapaz de aprender ningún patrón no lineal. Solo al introducir no linealidad la red puede aproximar funciones arbitrariamente complejas.
\end{warningbox}

\begin{figure}[H]
\centering
\includegraphics[width=0.85\textwidth]{slides-images/slide-028.jpg}
\caption{Las funciones de activación no lineales permiten que la red neuronal aprenda fronteras de decisión no lineales complejas\protect\footnotemark}
\end{figure}
\footnotetext{Diapositivas, página 28. Izquierda: la activación lineal solo puede producir fronteras de decisión lineales; derecha: la activación no lineal puede aproximar funciones arbitrariamente complejas.}

\subsection{Intuición geométrica del perceptrón}

Amini construye la intuición mediante un ejemplo bidimensional concreto. Considérese un perceptrón con dos entradas $(x_1, x_2)$, con parámetros $w_0=1$, $\mathbf{W}=[3, -2]^T$:

$$\hat{y} = g(1 + 3x_1 - 2x_2)$$

\begin{figure}[H]
\centering
\includegraphics[width=0.85\textwidth]{slides-images/slide-032.jpg}
\caption{Ejemplo de perceptrón: la frontera de decisión divide el espacio bidimensional en dos regiones\protect\footnotemark}
\end{figure}
\footnotetext{Diapositivas, página 32.}

La parte lineal dentro del paréntesis, $z = 1 + 3x_1 - 2x_2$, define una recta en el espacio bidimensional. Esta recta es la \textbf{frontera de decisión}:
\begin{itemize}
\item Cuando $z > 0$, después de la activación Sigmoid $\hat{y} > 0.5$, se predice la clase positiva
\item Cuando $z < 0$, $\hat{y} < 0.5$, se predice la clase negativa
\item Cuando $z = 0$, $\hat{y} = 0.5$, exactamente sobre la frontera de decisión
\end{itemize}

Tomando como ejemplo la entrada $\mathbf{X}=[-1, 2]^T$: $\hat{y} = g(1 + 3\times(-1) - 2\times 2) = g(-6) \approx 0.002$, muy por debajo de 0.5, por lo que este punto se clasifica como clase negativa.

\begin{knowledgebox}{Limitación de un solo perceptrón}
Un solo perceptrón solo puede aprender fronteras de decisión lineales (una recta en el espacio bidimensional, un hiperplano en espacios de mayor dimensión). Esto significa que no puede resolver problemas no separables linealmente como el XOR. Para tratar patrones más complejos, es necesario combinar múltiples perceptrones en una red de varias capas.
\end{knowledgebox}

\subsection{Resumen de la sección}

El perceptrón es la unidad mínima de la red neuronal; su propagación hacia adelante consta de tres pasos: suma ponderada, suma del sesgo y paso por la función de activación. El término de sesgo permite que la frontera de decisión no tenga que pasar por el origen, y la función de activación no lineal permite que la red aprenda patrones complejos. Un solo perceptrón solo puede representar un clasificador lineal.
